% !TeX spellcheck = en_US
\documentclass[french]{yLectureNote}

\title{Électromagnétisme 3}
\subtitle{Physique}
\author{Paulhenry Saux}
\date{\today}
\yLanguage{Français}
%jesse.groenen@cemes.fr
\professor{M.Brut}
\usepackage{graphicx}%----pour mettre des images
\usepackage[utf8]{inputenc}%---encodage
\usepackage{geometry}%---pour modifier les tailles et mettre a4paper
%\usepackage{awesomebox}%---pour les boites d'exercices, de pbq et de croquis ---d\'esactiv\'e pour les TP de PC
\usepackage{tikz}%---pour deiffner + d\'ependance de chemfig
% \usepackage{tabularx}%---pour dimensionner automatiquement les tableaux avec variable X
\usepackage{awesomebox}%---Pour les boites info, danger et autres
\usepackage{menukeys}%---Pour deiffner les touches de Calculatrice
\usepackage{fancyhdr}%---pour les en-t\^ete personnalis\'ees
\usepackage{blindtext}%---pour les liens
\usepackage{hyperref}%---pour les liens (\`a mettre en dernier)
\usepackage{caption}%---pour la francisation de la l\'egende table vers Tableau
\usepackage{pifont}
\usepackage{array}%---pour les tableaux
\usepackage{lipsum}
\usepackage{yFlatTable}
\usepackage{multicol}
\usepackage{tabularx}
\usepackage{booktabs}
\newcommand{\Lim}[1]{\lim\limits_{\substack{#1}}\:}
\renewcommand{\vec}{\overrightarrow}
\newcommand{\N}[0]{\mathbb{N}}
\newcommand{\dd}{\mathrm{d}}
\newcommand{\norm}[1]{||\vec{#1}||}
\newcommand{\fo}{\psi(\vec{r},t)}
\newcommand{\foe}{\psi(\vec{r},t)\*}
\newcommand{\HH}{\hat{H}}
\newcommand{\hb}{\hbar}
\newcommand{\lap}{\nabla^2}
\newcommand{\lapcc}{\frac{\partial^2 }{\partial x^2}+\frac{\partial^2 }{\partial y^2}+\frac{\partial^2 }{\partial z^2}}
\newcommand{\E}{\vec{E}(M)}
\newcommand{\B}{\vec{B}(M)}
\newcommand{\V}{V(M)}
\newcommand{\er}{\vec{e_{\rho}}}
\newcommand{\ep}{\vec{e_{\varphi}}}
\newcommand{\pip}{\pi^+}
\newcommand{\pim}{\pi^-}
\newcommand{\mv}{\mathcal{V}}
\DeclareMathOperator{\Div}{Div}
\newcommand{\rot}{\vec{\mathrm{rot}}}
\newcommand{\grad}{\vec{\mathrm{grad}}}
%\setcounter{chapter}{3}
\begin{document}
%\chapter{Puissance et énergie du champ Em}
\setcounter{chapter}{2}
\chapter{Milieux magnétiques}
\criticalInfo{Rappel}{Le moment magnétique est défini par la règle du tire-bouchon, il vaut $vec{\mu}=I\vec{S}$}
\section{Introduction}
\begin{definition}[Moment magnétique oribital]
    Déplacement de particule chargée dans l'espace. Le mouvement des électrons liés forment des courants localisés.
\end{definition}
\begin{definition}[Moment magnétique de spin]
    Moment magnétique intrinsqèque ses particules constitutives de l'atome. Le Spin est un DDL supplémentaire, quantique.
\end{definition}
Le moment magnétique des édifices est la somme des moments magnétiques orbitaux et de spin de leur constituants\marginTips{Il est essentiellement du aux électrons célibataires.}. 

\begin{definition}[Milieu diamagnétique]
    Il n'y a pas de moment magnétique permanent. Avec un champ magnétique appliqué, il y a un moment induit.
\end{definition}
\begin{definition}[Milieu paramagnétique]
    Il n'y a pas d'aimantation tant qu'on applique pas de champ magnétique. Dans ce cas, les moments sont selon le champ appliqué.
\end{definition}
\begin{definition}[Milieu ferro-magnétique]
    Il y a une aimantation m\^eme en l'absence de champ appliqué
\end{definition}
\section{Aimantation}
\subsection{Vecteur aimantation volumique}
\begin{definition}[Vecteur aimantation volumique]
    $$\delta \vec{\mu}=\vec{M}\delta v \iff \dd \vec{\mu}=\vec{M}\dd v$$ avec $\vec{M}$ le vecteur d'aimantation volumique, densité volumique de moment dipolaire magnétique. C'est une grandeur macroscopique.
\end{definition}
\subsection{Potentiel vecteur et champ magnétique crées par un milieu aimanté}
\begin{definition}[Potentiel vecteur $\vec{A_m}$]
    $$\vec{A_m}(\vec{r}) = \frac{\mu_0}{4\pi}\iiint_V\frac{\vec{M}(\vec{r'})\wedge (\vec{r}-\vec{r'})}{||\vec{r'}-\vec{r}||}\dd v'$$
\end{definition}
\begin{definition}[Densité volumique de courant]
    $$\vec{j_M} = \rot(\vec{M})$$ en volume
    $$\vec{j_{S,M}} = \vec{M}\wedge \vec{n_{ext}}$$ à la surface
\end{definition}
Il y a une équivalebce entre les potentiels vecteurs crées par un milieu aimanté avec l'aimantation volumique $\vec{M}$ et par les dictribution de courants d'aimantation.
\begin{theorem}[Relation locale $\vec{A_m}$ et champ magnétique]
    $$\vec{B_m} = \rot{\vec{A_m}}$$
\end{theorem}
\subsection{Interprétation}
Rappel : $i\dd l = j_s \dd S = j\dd V$ et $\vec{j} = \rho\vec{v}$
\begin{definition}[Densité surfacique de courant]
    $$j_{S,M}(P) = M(P)$$ et $$\vec{j_{s,M}}(N) = \vec{M}(N)\wedge \vec{n_{ext}}(N)$$
\end{definition}
\subsubsection{Dans un volume}
On fait la somme des petits courants. Si le courant est unforme, à l'inerface, ces courants sont opposés donc sont nuls. Il n'y a pas de courant en volume 
mais seulement surfacique.

Dans le cas non uniforme selon une direction. Dans la direction constante, on peut écrire :
$i_y = -ca \frac{\partial M_Z}{\partial x}(P)$

On a alors $\vec{j_M}(P) = -\frac{\partial M_Z}{\partial x}(P)\vec{e_y}$

On peut faire de m\^eme avec une variation selon l'autre axe.

On retrouve alors la formule précédente avec le rot.
\section{Excitation magnétique}
\subsection{Équation de maxwell dans un milieu magnétique (statique)}
On a 
\begin{itemize}
    \item $\Div \vec{B} = 0$
    \item $\rot{\vec{B}} = \mu_0 \vec{j_{total}}$
\end{itemize}
Il y a 2 termes : $\vec{j_{total}} = \vec{j_{libre}}+\vec{j_{M}}$

\begin{definition}[Excitation magnétique]
    $$\vec{H} = \frac{\vec{B}}{\mu_0}-\vec{M}$$
\end{definition}
On peut donc écrire 
\begin{theorem}
    $$\rot{\vec{H}} = \vec{j_{libre}}$$
\end{theorem}
\subsection{Théorème d'Ampère}
\begin{theorem}[Théorème d'Ampère]
   $$ \oint_C \vec{H}\cdot \vec{\dd l} = I_{libre, enlacé}$$
   I est algébrique, bien évidement
\end{theorem}
\subsection{Relation de passage}
\begin{theorem}[Relation de passage]
    $$\vec{n_{12}}\cdot(\vec{B_2}-\vec{B_1}) = 0$$
    $$\vec{n_{12}}\wedge(\vec{H_2}-\vec{H_1}) = \vec{j_{s,libre}}$$
\end{theorem}
\section{Milieux magnétiques, Linéaires, Homogènes, Isotropes}
\subsection{Susceptibilité $\chi_M$}
2 définitions co-existent :
\begin{itemize}
    \item $\vec{M} = \chi_M \vec{H}$ : Intéret : H est controlé expérimental
    \item $\vec{M} = \chi_m^*\frac{\vec{B}}{\mu_0}$
\end{itemize}
\subsection{Perméabilité}
\begin{definition}[Perméabilité]
    $$\vec{B} = \mu_0\mu_r \vec{H} = \mu\vec{H}$$ avec la Perméabilité, relative et absolue
\end{definition}
Comme il y a 2 définitions, il y a 2 liens :
\begin{definition}[Expression de $\mu_r$]
    $$\mu_r = 1+\chi_m = \frac{1}{-X_m^*}$$, donc $$\chi_m = \frac{\chi_m^*}{1-\chi_m^*}$$
\end{definition}
% \chapter{Polarisation milieux diélectriques : aspects microscopiques}

% Un conducteur parfait est équivalent à un milieu infiniment polarisable : $\chi_e \to \infty$. L'écrantage est total : $\vec{E_{int}} = \vec{0}$.
%  C'est équivalent en polarisation alors qu'il n'y a pas de dipôle élémentaire. C'est un réservoir illimité d'électrons libres.
% \section{Polarisabilité}
% \subsection{Définition}
% On se place en échelle microscopique
% \begin{definition}
%     $$\vec{p_j} = \epsilon_0 \alpha_j \vec{e_{loc}(M_j)}$$ avec $\vec{p_j}$ le moment dipolaire électrique induit sur l'élément j , $\vec{e_{loc}}$ le champ électrique micorcopique perçu au site $j$ par le constituant élémentaire
%     et $\alpha_j$ la Polarisabilité de l'élément j.
% \end{definition}
% \warningInfo{Champ local}{
%     $\vec{e_{loc}}$ correspond au champ appliqué $\vec{E_a}$ + champ microscopique crée par les constituants du milieu entourant $M_j$.

%     $\alpha_j$ est une constante positive.
% }
% \subsection{Bilan}
% \begin{itemize}
%     \item Polarisation électronique : Modification de la répartiiton des charges internes à haque atome ou ion, toujours présente
%     \item Polarisation atomique : Déplacement des atomes/ions par rapport à leur posiuton d'équilibre
%     \item Polarisation d'orienrttaion : Orientation du moment dipolaire permananr d'une molcéule polaire. Nécessite une liberté de mouvement. Cela concerne donc les fluides.
% \end{itemize}
\end{document}
